\begin{center}
\begin{tikzpicture}[font=\small,>=Stealth,line cap=round,line join=round]
\path[use as bounding box] (-6.5,-3.55) rectangle (6.5,3.55);
\draw[black!20,rounded corners=3pt] (-6.5,-3.55) rectangle (6.5,3.55);
\node[font=\small\bfseries] at (0,3.08) {完整直线：自身与每条垂线都是对称轴};

\begin{scope}[rotate=63.43495]
\draw[CExplanation,very thick] (-2.6,0)--(2.6,0);
\foreach \u in {-1.2,0,1.2} {\draw[CProof,dashed] (\u,-1.5)--(\u,1.5);}
\draw[CStatement,<->,thick] (-.9,0)--(.9,0);
\fill[CStatement] (-.9,0) circle(2pt); \fill[CStatement] (.9,0) circle(2pt);
\draw (0,.23)--(.23,.23)--(.23,0);
\end{scope}
\node[anchor=west] at (2.2,2.1) {蓝色实线：$y=2x$};
\node[anchor=west] at (2.2,1.35) {虚线：三条对称轴};
\node[anchor=west] at (2.2,.6) {这样的轴有无数条};

\node[align=center] at (0,-3.05) {关于自身，每个点不动；关于任一垂线，直线上的点成对交换。};
\end{tikzpicture}
\end{center}
